\chapter{Notation and conventions}
\label{book:appendixB}

Chapters define their symbols where they first use them. This appendix lists the
conventions that hold across chapters, so that a symbol borrowed from an earlier
chapter keeps its meaning.

\section{Sequences, positions and strands}
\begin{center}\small
\begin{tabular}{@{}P{0.24\linewidth}P{0.70\linewidth}@{}}
\toprule
Symbol & Meaning\\\midrule
$\Sigma$, $V$ & Token alphabet and its size; $\{A,C,G,T\}$ for canonical DNA\\
$x=(x_0,\ldots,x_{L-1})$ & A record of length $L$, indexed from zero\\
$[a,b)$ & Half-open base interval: $a\le p<b$\\
$c(\cdot)$, $R(x)$ & Base complement and reverse complement, $R(x)_i=c(x_{L-1-i})$\\
$u_t$, $y_t$ & Model input and next-token target at step $t$ ($y_t=x_{t+1}$ in causal mode)\\
$m_t$, $v_t$ & Validity or target masks; a mask is a weight or a legality flag, never a label\\
$\mathcal F(p)$ & Mutation footprint: the tokens whose spans contain base $p$\\
\bottomrule
\end{tabular}
\end{center}

\section{Tensors, models and state}
\begin{center}\small
\begin{tabular}{@{}P{0.24\linewidth}P{0.70\linewidth}@{}}
\toprule
Symbol & Meaning\\\midrule
$[B,T,D]$ & Batch, time (positions) and feature width; $H$ heads of width $d_h=D/H$\\
$\theta$, $\nabla_\theta\mathcal L$ & Learned parameters and the gradient of a scalar loss\\
$h_t$, $M_t$ & Recurrent fast state and memory loci after step $t$\\
$K^\ell,V^\ell$ & Layer-$\ell$ key and value cache, shape $[H,t,d_h]$\\
$(a,b)$, $\star$ / $\otimes$ & Affine pair $h\mapsto a\odot h+b$ and its chronological composition\\
$\odot$ & Elementwise product\\
$\sigma(\cdot)$ & Logistic sigmoid $1/(1+e^{-z})$\\
$\operatorname{softmax}$ & Row-wise $\exp(z_j)/\sum_k\exp(z_k)$ over \emph{allowed} entries\\
$\operatorname{LN}$ & Layer normalization over the feature axis of one token\\
\bottomrule
\end{tabular}
\end{center}
In mathematics vectors are columns. When the text follows PyTorch storage, batches are
rows and a linear layer applies the transpose of its stored weight matrix. Chapters say
which convention a formula uses.

\section{Probability, evaluation and units}
\begin{center}\small
\begin{tabular}{@{}P{0.24\linewidth}P{0.70\linewidth}@{}}
\toprule
Symbol & Meaning\\\midrule
$\mathcal L$, nats & Mean negative log-likelihood with natural logarithms\\
BPB & Bits per raw base: summed loss divided by $B\ln2$ for $B$ represented bases\\
$D_{\mathrm{KL}}(p\Vert q)$ & Kullback--Leibler divergence $\sum_xp(x)\log(p(x)/q(x))$\\
$\mathrm{TV}(p,q)$ & Total variation distance $\tfrac12\sum_x|p_x-q_x|$\\
$\widehat\Delta$ & An estimated paired difference between two predictors\\
$p_{\mathrm{upper}}=1-\alpha^{1/n}$ & One-sided upper bound on a failure rate after $n$ independent zero-failure trials\\
MAC, FLOP & Multiply--accumulate; floating-point operation (a MAC is often counted as two FLOPs)\\
KiB, MiB / KB, MB & Powers of 1024 / powers of 1000 bytes\\
tick & One logical step of a reference program, not a time measurement\\
\bottomrule
\end{tabular}
\end{center}

\section{Typographic conventions}
\texttt{Monospace} marks code, file names and literal sequences. A \Term{bold term} is
defined where it appears and is listed in the index. Listings are extracted from the
tested reference programs at build time and are never retyped. ``Generated'' tables
come from running the chapter's reference code; ``assigned'' numbers are chosen by hand
to make arithmetic visible and are not measurements.
